\documentclass[10pt,a4paper]{article}

\usepackage[a4paper,top=0.45cm,bottom=0.55cm,left=0.62cm,right=0.62cm]{geometry}
\usepackage[T1]{fontenc}
\usepackage[utf8]{inputenc}
\usepackage{lmodern}
\usepackage{multicol}
\usepackage{xcolor}
\usepackage{listings}
\usepackage[explicit]{titlesec}
\usepackage{enumitem}
\usepackage{tabularx}
\usepackage{booktabs}
\usepackage{array}
\usepackage{fancyhdr}
\usepackage{microtype}
\usepackage{needspace}
\usepackage{amsmath}

\definecolor{Navy}{HTML}{16324F}
\definecolor{Blue}{HTML}{22577A}
\definecolor{PaleBlue}{HTML}{EAF3F8}
\definecolor{CodeBg}{HTML}{F4F6F8}
\definecolor{CodeFrame}{HTML}{D9E0E5}
\definecolor{CodeGreen}{HTML}{357A38}
\definecolor{CodePurple}{HTML}{7A3E9D}
\definecolor{CodeGray}{HTML}{66727A}
\definecolor{Warn}{HTML}{8A3B12}

\pagestyle{fancy}
\fancyhf{}
\renewcommand{\headrulewidth}{0pt}
\fancyfoot[C]{\color{CodeGray}\fontsize{8}{9}\selectfont Python 3.12+ Quick Reference \textbullet\ \thepage}
\setlength{\footskip}{0.42cm}

\setlength{\columnsep}{0.42cm}
\setlength{\columnseprule}{0.15pt}
\def\columnseprulecolor{\color{CodeFrame}}
\setlength{\parindent}{0pt}
\setlength{\parskip}{0.9pt}
\setlist{nosep,leftmargin=1.15em,labelsep=.35em,topsep=1pt}
\raggedcolumns

\titleformat{\section}
  {\needspace{4\baselineskip}\bfseries\color{white}\fontsize{11.2}{12.2}\selectfont}
  {}{0pt}{\colorbox{Navy}{\parbox{\dimexpr\linewidth-2\fboxsep\relax}{\strut#1}}}
\titlespacing*{\section}{0pt}{4.5pt}{2.2pt}
\titleformat{\subsection}
  {\needspace{3\baselineskip}\bfseries\color{Blue}\fontsize{10}{11}\selectfont}
  {}{0pt}{#1}
\titlespacing*{\subsection}{0pt}{3pt}{1pt}

\lstdefinestyle{python}{
  language=Python,
  basicstyle=\ttfamily\fontsize{10}{10.7}\selectfont,
  keywordstyle=\bfseries\color{Blue},
  commentstyle=\itshape\color{CodeGreen},
  stringstyle=\color{CodePurple},
  showstringspaces=false,
  columns=fullflexible,
  keepspaces=true,
  breaklines=true,
  breakatwhitespace=true,
  upquote=true,
  backgroundcolor=\color{CodeBg},
  frame=single,
  rulecolor=\color{CodeFrame},
  framerule=0.25pt,
  framesep=1.5pt,
  xleftmargin=1pt,
  xrightmargin=1pt,
  aboveskip=1.0pt,
  belowskip=1.0pt,
  tabsize=4
}
\lstset{style=python}

\newcolumntype{L}{>{\raggedright\arraybackslash}X}
\newcolumntype{R}[1]{>{\raggedright\arraybackslash}p{#1}}
\newcommand{\py}[1]{\texttt{#1}}
\newcommand{\hint}[1]{\textcolor{CodeGray}{#1}}
\newcommand{\warn}[1]{\textcolor{Warn}{\bfseries #1}}
\newsavebox{\tblbox}
\newenvironment{tbl}{\par\begin{lrbox}{\tblbox}\begin{minipage}{\linewidth}\compacttable}%
  {\end{minipage}\end{lrbox}\noindent\raisebox{\dp\tblbox}{\usebox{\tblbox}}\par}
\newcommand{\compacttable}{\fontsize{10}{11.2}\selectfont\setlength{\tabcolsep}{2.5pt}\renewcommand{\arraystretch}{1.08}}

\begin{document}

\begin{center}
  {\color{Navy}\bfseries\fontsize{16}{17}\selectfont PYTHON 3.12+ CHEAT SHEET}\par
  \vspace{1pt}
  {\color{CodeGray}\fontsize{9}{10}\selectfont Practical syntax, idioms, and standard-library patterns for coding interviews}
\end{center}
\vspace{-3pt}

\begin{multicols}{2}

\section{1. Python Basics}

\begin{lstlisting}
# comment; names bind to objects
x, name, active = 42, "Ada", True
a = b = 0                 # multiple assignment
missing = None            # absence of a value
\end{lstlisting}

\begin{tbl}
\begin{tabularx}{\linewidth}{@{}R{.28\linewidth}L@{}}
\toprule
\textbf{Type} & \textbf{Typical literal / conversion} \\
\midrule
\py{int}, \py{float}, \py{complex} & \py{12}, \py{3.5}, \py{2+3j}; \py{int("12")}, \py{float(x)} \\
\py{bool}, \py{NoneType} & \py{True}, \py{False}, \py{None}; \py{bool(x)} \\
\py{str}, \py{bytes} & \py{"text"}, \py{b"raw"}; \py{str(x)}, \py{bytes(s,"utf-8")} \\
\py{list}, \py{tuple} & \py{[1,2]}, \py{(1,2)}; mutable / immutable sequence \\
\py{dict}, \py{set} & \py{\{"a":1\}}, \py{\{1,2\}}; mapping / unique values \\
\bottomrule
\end{tabularx}
\end{tbl}

\begin{lstlisting}
type(x) is int                 # exact type
isinstance(x, (int, float))    # type or subclass
a == b                         # equal values
a is b                         # identical object (use: x is None)
\end{lstlisting}

\textbf{Falsy:} \py{False}, \py{None}, numeric zero, and empty strings/containers. Everything else is truthy.

\subsection{References, copies, and hashability}

Names refer to objects; assignment binds another name, not a copy. \py{==} compares values; \py{is} compares identity; \py{id(obj)} identifies an object during its lifetime.

\begin{lstlisting}
a = [1, 2]; b = a
b.append(3)                  # a is now [1, 2, 3]
import copy
shallow = a.copy()           # also a[:] or copy.copy(a)
deep = copy.deepcopy(a)      # recursively copies nested objects
\end{lstlisting}

Shallow copies share nested references. Mutable: \py{list}, \py{dict}, \py{set}. Immutable: \py{int}, \py{float}, \py{bool}, \py{str}, \py{tuple}, \py{frozenset} (a tuple can still contain mutable objects).

\textbf{Hashable keys/elements:} dictionaries and sets use hashing. Usually hashable: numbers, strings, \py{frozenset}, and tuples of hashable values. Lists, dicts, and sets are unhashable. \py{hash(x)} gets a hash value.

\begin{lstlisting}
d = {(1, 2): "valid"}       # tuple key
# {[1, 2]: "invalid"}      # TypeError: unhashable list
\end{lstlisting}

\section{2. Strings}

Strings are immutable Unicode sequences; operations return new strings.

\begin{lstlisting}
s = "Python"
s[0], s[-1]        # ('P', 'n')
s[1:4], s[::-1]    # ('yth', 'nohtyP')
len(s), "th" in s  # (6, True)
\end{lstlisting}

\begin{tbl}
\begin{tabularx}{\linewidth}{@{}R{.46\linewidth}L@{}}
\toprule
\textbf{Expression} & \textbf{Result / purpose} \\
\midrule
\py{s.split(",")} & split into a list \\
\py{"-".join(parts)} & join strings with separator \\
\py{s.strip()} & trim surrounding whitespace \\
\py{s.replace("a","b")} & replace all occurrences \\
\py{s.find("th")} & first index (\py{2}), or \py{-1} \\
\py{s.startswith("Py")} & prefix test \\
\py{s.endswith(".py")} & suffix test \\
\py{s.upper()}, \py{s.lower()} & change case \\
\py{s.capitalize()}, \py{s.title()} & first / each word capitalized \\
\py{s.count("a")} & count non-overlapping matches \\
\py{s.isalpha()}, \py{s.isdigit()} & letters only / digits only \\
\py{s.isalnum()}, \py{s.isspace()} & letters or digits / whitespace only \\
\py{"x" in s} & substring membership \\
\bottomrule
\end{tabularx}
\end{tbl}

\begin{lstlisting}
name, age = "Alice", 30
f"{name} is {age} years old"
f"{1234.5:,.2f}"       # '1,234.50'
f"{0.875:.1%}"         # '87.5%'
\end{lstlisting}

Alignment and display: \py{f"\{x:>10\}"} (right), \py{f"\{x:<10\}"} (left), \py{f"\{x:\textasciicircum{}10\}"} (center), \py{f"\{n:08d\}"} (zero-pad), \py{f"\{x=!r\}"} (debug).

\begin{lstlisting}
message = """A multiline string
may span several lines."""
raw_path = r"C:\new\data"    # backslashes stay literal
\end{lstlisting}

\py{ord("a") == 97}; \py{chr(97) == "a"}. Convert characters to/from code points for letter arithmetic.

\section{3. Numbers, Math, and Bits}

\begin{tbl}
\begin{tabularx}{\linewidth}{@{}R{.42\linewidth}L@{}}
\toprule
\textbf{Syntax} & \textbf{Meaning} \\
\midrule
\py{a+b}, \py{a-b}, \py{a*b}, \py{a/b} & arithmetic; \py{/} returns float \\
\py{a//b}, \py{a\%b}, \py{a**b} & floor division, remainder, power \\
\py{round(x, 2)}, \py{abs(x)} & round; magnitude \\
\py{min(xs)}, \py{max(xs)}, \py{sum(xs)} & aggregate values \\
\bottomrule
\end{tabularx}
\end{tbl}

\begin{lstlisting}
import math
math.sqrt(9); math.ceil(2.1); math.floor(2.9)
math.isclose(a, b, rel_tol=1e-9)
math.pi, math.e, math.log(x), math.sin(angle)
\end{lstlisting}

\py{//} and \py{\%} round toward negative infinity: \py{-7 // 2} is \py{-4} and \py{-7 \% 2} is \py{1}. \py{divmod(a, b)} returns both at once.

\subsection{Bitwise operators}

These act on the binary bits of an \py{int}. Python integers are signed and unbounded, so they behave like two's complement with infinite sign extension.

\begin{tbl}
\begin{tabularx}{\linewidth}{@{}R{.20\linewidth}L@{}}
\toprule
\textbf{Syntax} & \textbf{Meaning} \\
\midrule
\py{a \& b} & AND: 1 where both bits are 1 \\
\py{a | b} & OR: 1 where either bit is 1 \\
\py{a \textasciicircum{} b} & XOR: 1 where the bits differ \\
\py{\textasciitilde{}a} & NOT: invert every bit; equals \py{-a - 1} \\
\py{a <{}< n} & shift left; same as \py{a * 2**n} \\
\py{a >{}> n} & shift right; same as \py{a // 2**n} \\
\bottomrule
\end{tabularx}
\end{tbl}

\begin{lstlisting}
x = 0b1010            # 10; also 0o12, 0xA, 1_000
bin(10), hex(255)     # ('0b1010', '0xff')
int("1010", 2)        # 10
10 & 6, 10 | 6, 10 ^ 6      # (2, 14, 12)
~10, 10 << 1, 10 >> 1       # (-11, 20, 5)
(10).bit_length(), (10).bit_count()   # (4, 2)
\end{lstlisting}

\begin{lstlisting}
mask = 1 << i        # a 1 in position i only
n | mask             # set bit i
n & ~mask            # clear bit i
n ^ mask             # toggle bit i
(n >> i) & 1         # read bit i -> 0 or 1
n & (n - 1)          # clear the lowest set bit
n & -n               # isolate the lowest set bit
n > 0 and n & (n - 1) == 0    # exact power of two?
\end{lstlisting}

Bitwise operators bind tighter than comparisons but looser than arithmetic, so \py{n \& 1 == 0} groups as \py{(n \& 1) == 0}. \py{>{}>} preserves the sign and there is no unsigned \py{>{}>{}>}; mask with \py{\& 0xFF} to emulate a fixed width. For named flags, prefer \py{enum.IntFlag} over bare constants.

\section{4. Lists}

Mutable ordered sequences; indexing and slicing match strings.

\begin{lstlisting}
items = ["a", "b", "c"]
items[0], items[-1], items[1:]
items[1] = "B"; "a" in items
empty = []; shallow = items.copy()   # see copying in Basics
\end{lstlisting}

\begin{tbl}
\begin{tabularx}{\linewidth}{@{}R{.47\linewidth}L@{}}
\toprule
\textbf{Mutation} & \textbf{Effect} \\
\midrule
\py{xs.append(x)} & add one item at end \\
\py{xs.extend(more)} & append all items \\
\py{xs.insert(i,x)} & insert before index \py{i} \\
\py{xs.remove(x)} & delete first \py{x}; raises if absent \\
\py{x = xs.pop(i)} & delete and return item (last by default) \\
\py{xs.clear()} & remove every item \\
\py{xs.reverse()} & reverse in place \\
\py{xs.sort(key=..., reverse=True)} & sort in place \\
\bottomrule
\end{tabularx}
\end{tbl}

\begin{lstlisting}
nums = [4, 1, 3, 2]
ordered = sorted(nums)          # new list
nums.sort(reverse=True)         # mutate nums
\end{lstlisting}

Useful slices: \py{xs[:n]} first $n$; \py{xs[n:]} rest; \py{xs[::2]} every other; \py{xs[::-1]} reversed copy; \py{del xs[i:j]} delete range.

\section{5. Tuples and Sets}

\subsection{Tuples}
Immutable sequences are ideal for fixed records and unpacking. A one-item tuple needs a comma.

\begin{lstlisting}
point = (10, 20); single = (10,)
x, y = point
first, *middle, last = values
def bounds(xs): return min(xs), max(xs)
\end{lstlisting}

\subsection{Sets}
Mutable collections of unique hashable values; no positional indexing.

\begin{lstlisting}
a, b = {1, 2, 3}, {3, 4}
a.add(5); a.discard(9)    # no error if absent
a.remove(2)               # KeyError if absent
a | b, a & b, a - b       # union, intersection, difference
\end{lstlisting}

\py{x in a} tests membership; \py{set(items)} removes duplicates (order not preserved). Use \py{frozenset(...)} for an immutable, hashable set.

\section{6. Dictionaries}

Mappings preserve insertion order and require hashable keys.

\begin{lstlisting}
user = {"name": "Alice", "age": 30}
name = user["name"]             # KeyError if absent
age = user.get("age", 0)        # safe default
user["active"] = True           # add / update
user.update(age=31, city="Paris")
\end{lstlisting}

\begin{tbl}
\begin{tabularx}{\linewidth}{@{}R{.45\linewidth}L@{}}
\toprule
\textbf{Expression} & \textbf{Purpose} \\
\midrule
\py{k in d} & key membership \\
\py{d.keys()}, \py{d.values()} & dynamic views \\
\py{d.items()} & \py{(key, value)} pairs \\
\py{d.pop(k, default)} & remove key and return value \\
\py{del d[k]} & remove key; raises if absent \\
\py{d.setdefault(k, v)} & get existing or insert default \\
\py{d1 | d2} & new merged dict; right side wins \\
\bottomrule
\end{tabularx}
\end{tbl}

\begin{lstlisting}
for key, value in user.items():
    print(key, value)
lengths = {word: len(word) for word in words}
merged = defaults | overrides
\end{lstlisting}

\subsection{Counting and grouped values}

\begin{lstlisting}
from collections import Counter, defaultdict
counts = Counter(words)          # counts.most_common(3)
groups = defaultdict(list)
for item in rows: groups[item.key].append(item)
\end{lstlisting}

\section{7. Conditions}

\begin{lstlisting}
if score >= 90:
    grade = "A"
elif score >= 80:
    grade = "B"
else:
    grade = "C"
\end{lstlisting}

Comparisons: \py{== != < <= > >=}; logic: \py{and or not}. Comparisons chain naturally: \py{0 <= x < 10}. Short-circuit logic can supply defaults: \py{name = supplied or "Guest"}.

\begin{lstlisting}
status = "adult" if age >= 18 else "minor"
if ready and not cancelled:
    run()
\end{lstlisting}

\section{8. Loops}

\begin{lstlisting}
for item in items:
    process(item)
for i in range(0, 10, 2):       # 0, 2, ..., 8
    print(i)
\end{lstlisting}

\begin{lstlisting}
for i, value in enumerate(items, start=1):
    print(i, value)
for name, score in zip(names, scores, strict=True):
    print(name, score)
\end{lstlisting}

\begin{lstlisting}
for key, value in data.items():
    if value is None: continue
    if key == target: break
else:
    print("no break / not found")
\end{lstlisting}

\begin{lstlisting}
attempts = 3
while attempts > 0:
    if try_once(): break
    attempts -= 1
\end{lstlisting}

\py{break} exits the nearest loop; \py{continue} skips to its next iteration. Loop \py{else} runs only when no \py{break} occurs.

\section{9. Functions}

\begin{lstlisting}
def greet(name: str, greeting: str = "Hello") -> str:
    """Return a short greeting."""
    return f"{greeting}, {name}!"

greet("Ada", greeting="Hi")
\end{lstlisting}

\textbf{Parameters} may be positional or keyword. \py{/} marks positional-only; \py{*} marks following parameters keyword-only.

\begin{lstlisting}
def clamp(x, /, low=0, high=100, *, strict=False):
    if strict and not low <= x <= high:
        raise ValueError(x)
    return max(low, min(x, high))

def collect(*args, **kwargs):
    return args, kwargs           # tuple, dict
\end{lstlisting}

\begin{lstlisting}
nums = [3, 7]
range(*nums)                    # unpack positional args
opts = {"sep": " - ", "end": "!\n"}
print("a", "b", **opts)         # unpack keyword args
\end{lstlisting}

Lambdas are short expressions, best used as keys/callbacks: \py{sorted(users, key=lambda u: u["age"])}.

\warn{Mutable default trap:} defaults are created once.

\begin{lstlisting}
def add(x, items=None):
    items = [] if items is None else items
    items.append(x)
    return items
\end{lstlisting}

\subsection{Scope: LEGB, closures, decorators}

Name lookup checks \textbf{L}ocal, \textbf{E}nclosing, \textbf{G}lobal, then \textbf{B}uilt-in. \py{global} rebinds a module name; \py{nonlocal} rebinds a name in the nearest enclosing function scope.

\begin{lstlisting}
x = 10
def change_global():
    global x
    x += 1

def outer():
    y = 20
    def inner():
        nonlocal y
        y += 1
    inner()
    return y
\end{lstlisting}

Closures capture a variable, not its value when the function is created (late binding):

\begin{lstlisting}
funcs = []
for i in range(3):
    funcs.append(lambda: i)
[f() for f in funcs]         # [2, 2, 2]
# Fix inside loop: funcs.append(lambda i=i: i)
\end{lstlisting}

\begin{lstlisting}
from functools import wraps
def decorator(func):
    @wraps(func)
    def wrapper(*args, **kwargs):
        print("before")
        return func(*args, **kwargs)
    return wrapper

@decorator
def hello(): print("hello")
\end{lstlisting}

\py{@decorator} above \py{def f(): ...} means \py{f = decorator(f)}. \py{wraps} preserves the wrapped function's metadata.

\section{10. Comprehensions and Generators}

\begin{tbl}
\begin{tabularx}{\linewidth}{@{}R{.32\linewidth}L@{}}
\toprule
\textbf{Result} & \textbf{Pattern} \\
\midrule
list & \py{[f(x) for x in xs if test(x)]} \\
set & \py{\{f(x) for x in xs\}} \\
dict & \py{\{k: f(k) for k in keys\}} \\
generator & \py{(f(x) for x in xs)} (lazy) \\
\bottomrule
\end{tabularx}
\end{tbl}

\begin{lstlisting}
even_squares = [x*x for x in nums if x % 2 == 0]
by_id = {u.id: u for u in users}
total = sum(x*x for x in nums)  # no intermediate list
flat = [x for row in matrix for x in row]
\end{lstlisting}

Prefer a normal loop when the expression needs side effects or becomes hard to scan.

\section{11. Most Useful Built-ins}

\begin{tbl}
\begin{tabularx}{\linewidth}{@{}R{.40\linewidth}L@{}}
\toprule
\textbf{Purpose / calls} & \textbf{Use} \\
\midrule
\textbf{Inspect:} \py{type(x)}, \py{isinstance(x,T)} & exact type / subclass-aware test \\
\py{id(x)}, \py{hash(x)}, \py{len(x)} & identity / hash / size \\
\textbf{Convert:} \py{int(x)}, \py{float(x)} & number conversions \\
\py{str(x)}, \py{bool(x)} & text / truth conversion \\
\py{list(x)}, \py{tuple(x)} & sequence conversion \\
\py{set(x)}, \py{dict(x)} & unique values / mapping \\
\py{ord(c)}, \py{chr(n)} & character / code point \\
\textbf{Iterate:} \py{range(n)}, \py{iter(x)} & lazy range / iterator \\
\py{next(it, default)} & next value or default \\
\py{enumerate(xs)}, \py{zip(a,b)} & index pairs / paired items \\
\py{reversed(xs)}, \py{sorted(xs)} & reverse iterator / sorted list \\
\textbf{Aggregate:} \py{any(xs)}, \py{all(xs)} & some / every item truthy \\
\py{sum(xs)}, \py{min(xs)}, \py{max(xs)} & total / extrema; extrema take \py{key=} \\
\py{divmod(a,b)} & quotient and remainder \\
\bottomrule
\end{tabularx}
\end{tbl}

\begin{lstlisting}
oldest = max(users, key=lambda u: u.age)
has_errors = any(row.error for row in rows)
all_valid = all(validate(x) for x in items)
first = next((x for x in xs if test(x)), None)
\end{lstlisting}

\py{map(f, xs)} and \py{filter(test, xs)} return iterators; comprehensions are usually clearer.

\section{12. Exceptions}

Catch only errors you can handle; prefer specific exception types.

\begin{lstlisting}
try:
    value = int(text)
except ValueError as exc:
    value = 0
else:
    use(value)               # conversion succeeded
finally:
    cleanup()                # always runs
\end{lstlisting}

\begin{lstlisting}
if amount < 0:
    raise ValueError("amount must be nonnegative")
\end{lstlisting}

Multiple types: \py{except (TypeError, ValueError):}. Re-raise the current exception with bare \py{raise}; preserve cause with \py{raise NewError(...) from exc}.

\section{13. Files and pathlib}

Context managers close files even when an error occurs. Specify text encoding.

\begin{lstlisting}
with open("data.txt", "r", encoding="utf-8") as f:
    text = f.read()           # whole file
with open("out.txt", "w", encoding="utf-8") as f:
    f.write(text)
\end{lstlisting}

\py{"a"} appends; \py{"b"} selects binary data (for example \py{"rb"}). Iterate lines without loading the whole file:

\begin{lstlisting}
with open("data.txt", encoding="utf-8") as f:
    for line in f:
        process(line.rstrip("\n"))
\end{lstlisting}

\subsection{Custom context managers}

\py{with} calls \py{\_\_enter\_\_()} on entry and \py{\_\_exit\_\_()} on exit, even after an exception. Return \py{self} from \py{\_\_enter\_\_} for the \py{as} target.

\begin{lstlisting}
class Resource:
    def __enter__(self):
        print("acquire")
        return self
    def __exit__(self, exc_type, exc_value, traceback):
        print("release")

with Resource() as r:
    pass
\end{lstlisting}

\subsection{Prefer pathlib for paths}

\begin{lstlisting}
from pathlib import Path
p = Path("data") / "input.txt"
text = p.read_text(encoding="utf-8")
Path("out.txt").write_text("Hello", encoding="utf-8")
\end{lstlisting}

\begin{lstlisting}
p.exists(); p.is_file(); p.name; p.suffix; p.parent
p.mkdir(parents=True, exist_ok=True)
for file in Path(".").glob("*.py"):
    print(file)
\end{lstlisting}

\py{p.rglob("*.csv")} searches recursively. \py{p.rename(new)} moves a file; \py{p.unlink(missing\_ok=True)} deletes it if present.

\section{14. Modules and Environments}

\begin{lstlisting}
import statistics
import numpy as np
from pathlib import Path
from package.module import useful_name
\end{lstlisting}

Imports run module top-level code once and bind names. Put script entry behavior behind a guard:

\begin{lstlisting}
def main() -> None:
    ...

if __name__ == "__main__":
    main()
\end{lstlisting}

\subsection{Packages and \texttt{\_\_init\_\_.py}}

A directory with \py{\_\_init\_\_.py} is a regular import package; the file runs on first import and can re-export names. Namespace packages may omit it.

\begin{lstlisting}
# app/__init__.py
from .utils import helper      # re-export
# elsewhere
from app import helper
\end{lstlisting}

\subsection{Configuration with \texttt{.env}}

\py{.env} holds \py{KEY=value} settings; Python does not load it automatically. Use third-party \py{python-dotenv}:

\begin{lstlisting}
import os
from dotenv import load_dotenv
load_dotenv()
api_url = os.getenv("API_URL")   # None if unset
\end{lstlisting}

Existing environment values take precedence by default. Keep secret \py{.env} files out of version control; share a non-secret \py{.env.example} template.

\subsection{Virtual environments}

\py{venv} isolates each project's interpreter and installed packages. Run the commands for your platform:

\begin{lstlisting}
python3 -m venv .venv           # macOS/Linux
py -m venv .venv                # Windows
source .venv/bin/activate      # macOS/Linux
.venv\Scripts\Activate.ps1     # PowerShell
python -m pip install requests python-dotenv
deactivate
\end{lstlisting}

Exclude \py{.venv} from version control; activate it before installing or running project code.

\section{15. Classes and Data Classes}

\begin{lstlisting}
class Account:
    def __init__(self, owner: str, balance=0.0):
        self.owner, self.balance = owner, balance

    def deposit(self, amount: float) -> None:
        self.balance += amount
\end{lstlisting}

\begin{lstlisting}
class Savings(Account):
    def __init__(self, owner, rate):
        super().__init__(owner)
        self.rate = rate

    @property
    def interest(self): return self.balance * self.rate
\end{lstlisting}

\subsection{Properties: getter and setter}

\py{@property} runs a getter on attribute access; \py{@age.setter} runs a setter on assignment. Store the value under a different name to avoid recursion.

\begin{lstlisting}
class Person:
    def __init__(self, age): self.age = age
    @property
    def age(self): return self._age
    @age.setter
    def age(self, value):
        if value < 0: raise ValueError("age")
        self._age = value

p = Person(20)
p.age = 21; print(p.age)       # set, then get
\end{lstlisting}

\subsection{Is-a vs. has-a}

\py{Savings} \textbf{is an} \py{Account}: inheritance reuses or extends behavior. \py{Car} \textbf{has an} \py{Engine}: composition stores another object as an attribute.

\begin{lstlisting}
class Engine:
    def start(self): ...

class Car:
    def __init__(self):
        self.engine = Engine()    # has-a

car = Car()
car.engine.start()
\end{lstlisting}

\subsection{Polymorphism}

Different subclasses can respond to the same method call in their own way:

\begin{lstlisting}
class Animal:
    def speak(self): return "sound"
class Dog(Animal):
    def speak(self): return "woof"

for animal in (Animal(), Dog()):
    print(animal.speak())          # sound, woof
\end{lstlisting}

\subsection{Creation, class state, and methods}

\py{MyClass()} $\rightarrow$ \py{\_\_new\_\_(cls)} $\rightarrow$ new object $\rightarrow$ \py{\_\_init\_\_(self)}. \py{cls} is the class; \py{self} is the instance. \py{\_\_new\_\_} creates/returns an instance; \py{\_\_init\_\_} initializes the already-created instance.

\begin{lstlisting}
class Test:
    a = 100                       # class variable
    def __init__(self):
        self.b = 200              # instance variable

x = Test(); y = Test()
x.c = 300                         # dynamic instance attribute
\end{lstlisting}

\py{Test.a} is shared unless shadowed by \py{x.a = ...}; \py{x.b} and \py{y.b} belong to separate objects. Define \py{\_\_slots\_\_} to restrict allowed instance attributes.

\begin{lstlisting}
class MyClass:
    def method(self): ...        # receives instance
    @classmethod
    def create(cls): return cls() # receives class
    @staticmethod
    def helper(x): return x * 2  # receives neither
\end{lstlisting}

\subsection{Important dunder methods}

\begin{tbl}
\begin{tabularx}{\linewidth}{@{}R{.39\linewidth}L@{}}
\toprule
\textbf{Method} & \textbf{Triggered by / role} \\
\midrule
\py{\_\_new\_\_}, \py{\_\_init\_\_} & \py{C()} creates, then initializes \\
\py{\_\_str\_\_} & \py{print(obj)}, \py{str(obj)} \\
\py{\_\_repr\_\_} & \py{repr(obj)}, interactive display \\
\py{\_\_eq\_\_} & \py{a == b} \\
\py{\_\_hash\_\_} & \py{hash(obj)}, dict/set lookup \\
\py{\_\_len\_\_} & \py{len(obj)} \\
\py{\_\_getitem\_\_} & \py{obj[i]} \\
\py{\_\_setitem\_\_} & \py{obj[i] = x} \\
\py{\_\_iter\_\_} & \py{iter(obj)}, \py{for} \\
\py{\_\_next\_\_} & \py{next(obj)} \\
\py{\_\_call\_\_} & \py{obj()} \\
\py{\_\_enter\_\_}, \py{\_\_exit\_\_} & \py{with obj as x:} \\
\bottomrule
\end{tabularx}
\end{tbl}

\subsection{Singleton (interview pattern)}

\begin{lstlisting}
class Singleton:
    _instance = None
    def __new__(cls):
        if cls._instance is None:
            cls._instance = super().__new__(cls)
        return cls._instance

a = Singleton(); b = Singleton()
a is b                           # True
\end{lstlisting}

\py{\_\_new\_\_} returns the same object on later calls; \py{\_\_init\_\_} still runs each time unless guarded. A module often supplies shared state without a Singleton class.

Use data classes for data containers; they generate \py{\_\_init\_\_}, equality, and a useful representation.

\begin{lstlisting}
from dataclasses import dataclass, field

@dataclass(slots=True)
class User:
    name: str
    tags: list[str] = field(default_factory=list)
\end{lstlisting}

\section{16. Iterators and Generators}

An \emph{iterable} produces an iterator with \py{iter(obj)}; \emph{iterator} returns itself from \py{\_\_iter\_\_} and supplies values with \py{\_\_next\_\_}. Flow: iterable $\rightarrow$ \py{iter(obj)} $\rightarrow$ iterator $\rightarrow$ \py{next(it)} $\rightarrow$ value (or \py{StopIteration}).

\begin{lstlisting}
it = iter([10, 20])
next(it), next(it)       # 10, 20
next(it, None)           # None (default avoids StopIteration)
\end{lstlisting}

\begin{lstlisting}
class Counter:
    def __init__(self, n):
        self.n, self.i = n, 0
    def __iter__(self):
        return self
    def __next__(self):
        if self.i >= self.n:
            raise StopIteration
        self.i += 1
        return self.i
\end{lstlisting}

Generators produce lazy streams and retain local state:

\begin{lstlisting}
def countdown(n):
    while n:
        yield n
        n -= 1

squares = (x*x for x in range(1_000_000))
\end{lstlisting}

Generators are single-use; materialize with \py{list(gen)} only when needed.

\section{17. Standard Library Essentials}

\subsection{collections}

\begin{lstlisting}
from collections import Counter, defaultdict, deque
Counter("banana").most_common(2)  # [('a', 3), ('n', 2)]
groups = defaultdict(list)
queue = deque([1, 2]); queue.append(3); queue.popleft()
\end{lstlisting}

\subsection{heapq (min-heap / priority queue)}

A heap is an ordinary list kept in heap order, so \py{h[0]} is always the smallest item. Push and pop cost $O(\log n)$ and \py{heapify} costs $O(n)$; the list as a whole is \emph{not} sorted.

\begin{lstlisting}
import heapq
h = [5, 1, 3]
heapq.heapify(h)          # -> [1, 5, 3], rearranged in place
heapq.heappush(h, 2)
smallest = heapq.heappop(h)   # 1; peek instead with h[0]
heapq.heappushpop(h, 7)   # push then pop, in one pass
heapq.heapreplace(h, 7)   # pop then push; heap must be non-empty
\end{lstlisting}

\begin{lstlisting}
heapq.nsmallest(3, xs)                 # sorted; cheap for small k
heapq.nlargest(3, words, key=len)
list(heapq.merge(sorted_a, sorted_b))  # lazy merge of sorted inputs
\end{lstlisting}

\textbf{Python 3.12--3.13:} negate numbers (or a sort key) for a max-heap:

\begin{lstlisting}
maxh = [-x for x in xs]; heapq.heapify(maxh)
largest = -heapq.heappop(maxh)
\end{lstlisting}

\textbf{Python 3.14+:} \py{heapq} also has \py{heapify\_max}, \py{heappush\_max}, and \py{heappop\_max}.

Entries are compared as whole tuples, so insert a unique tie-breaker ahead of any payload that is not orderable:

\begin{lstlisting}
from itertools import count
tie = count()                 # unique, increasing
heapq.heappush(pq, (priority, next(tie), task))
priority, _, task = heapq.heappop(pq)   # lowest priority first
\end{lstlisting}

Use \py{queue.PriorityQueue} instead when several threads share the queue.

\subsection{datetime}

\begin{lstlisting}
from datetime import datetime, date, timedelta, timezone
now = datetime.now(timezone.utc); today = date.today()
dt = datetime.strptime("2026-08-08", "%Y-%m-%d")
dt.strftime("%d %b %Y"); tomorrow = today + timedelta(days=1)
\end{lstlisting}

Use timezone-aware datetimes for stored/shared timestamps.

\subsection{json}

\begin{lstlisting}
import json
data = json.loads('{"ok": true}')
text = json.dumps(data, indent=2, ensure_ascii=False)
data = json.loads(Path("data.json").read_text())
Path("out.json").write_text(json.dumps(data, indent=2))
\end{lstlisting}

\subsection{random}

\begin{lstlisting}
import random
n = random.randint(1, 6)       # inclusive endpoints
item = random.choice(items)
random.shuffle(items)          # in place
\end{lstlisting}

Use \py{secrets} rather than \py{random} for passwords, tokens, or security.

\subsection{Regular expressions}

Use raw strings for patterns; compile a pattern when reusing it often.

\begin{lstlisting}
import re
m = re.search(r"\b\d+\b", text)     # first match / None
nums = re.findall(r"\d+", text)     # all matches
clean = re.sub(r"\s+", " ", text).strip()
\end{lstlisting}

\subsection{itertools (compact combinatorics)}

\begin{lstlisting}
from itertools import chain, product, combinations
list(chain(xs, ys))
product("AB", repeat=2); combinations(items, 2)
\end{lstlisting}

\subsection{functools (memoization and helpers)}

\begin{lstlisting}
from functools import lru_cache, cache, wraps, partial, reduce

@cache                       # or @lru_cache(None)
def dp(i):
    ...                      # recursive subproblem
\end{lstlisting}

Memoization turns repeated recursive subproblems into cached lookups; use hashable arguments. \py{wraps} preserves decorator metadata; \py{partial(f, x)} fixes arguments; \py{reduce(f, xs)} folds a sequence.

\subsection{bisect (sorted-list search)}

\begin{lstlisting}
import bisect
nums = [1, 3, 5, 7]
bisect.bisect_left(nums, 5)     # 2: before equals
bisect.bisect_right(nums, 5)    # 3: after equals
bisect.insort(nums, 4)          # keep list sorted
\end{lstlisting}

Search position: $O(\log n)$; inserting into a list still costs $O(n)$.

\section{18. Working with Data and APIs}

\subsection{Web communication}

Client $\rightarrow$ HTTP request (method, URL, headers, optional body) $\rightarrow$ server $\rightarrow$ response (status, headers, body). \py{GET} reads; \py{POST} submits. Status: \py{2xx} success, \py{4xx} client error, \py{5xx} server error.

\subsection{HTTP requests and JSON}

\py{requests} is a third-party package; \py{json=} encodes a request body. The \py{json.loads/dumps} examples above use the standard library.

\begin{lstlisting}
import requests
url = "https://api.example.com/items"
params = {"q": "py"}
r = requests.get(url, params=params, timeout=5)
r.raise_for_status()          # fail on 4xx / 5xx
items = r.json()              # decode JSON
payload = {"name": "Ada"}
r = requests.post(url, json=payload, timeout=5)
r.raise_for_status()
\end{lstlisting}

\subsection{Network errors and resilience}

Exception classes below are in \py{requests.exceptions}.

\begin{tbl}
\begin{tabularx}{\linewidth}{@{}R{.44\linewidth}L@{}}
\toprule
\textbf{Exception} & \textbf{Meaning} \\
\midrule
\py{Timeout} & response took too long \\
\py{ConnectionError} & connection / DNS failure \\
\py{HTTPError} & bad status via \py{raise\_for\_status()} \\
\py{JSONDecodeError} & invalid JSON response \\
\bottomrule
\end{tabularx}
\end{tbl}

Set a timeout. Retry idempotent requests after timeouts, \py{429}, or transient \py{5xx} with bounded backoff; respect \py{Retry-After}. Do not blindly retry \py{POST}, which may duplicate a write.

\subsection{Storing data with SQLite}

\begin{lstlisting}
import sqlite3
con = sqlite3.connect("app.db")
schema = "CREATE TABLE IF NOT EXISTS t(name)"
try:
    with con:                 # commit or rollback
        con.execute(schema)
        sql = "INSERT INTO t VALUES (?)"
        con.execute(sql, ("Ada",))
        cur = con.execute("SELECT name FROM t")
        row = cur.fetchone()
finally:
    con.close()
\end{lstlisting}

Use \py{?} placeholders for data values, never string formatting in SQL. \py{fetchone()} returns one row or \py{None}; \py{fetchall()} returns all rows. \py{with con} does not close it.

\section{19. Common Python Patterns}

\begin{lstlisting}
a, b = b, a                         # swap
unique = list(dict.fromkeys(items)) # dedupe, keep order
\end{lstlisting}

\subsection{Stack and queue}

\begin{lstlisting}
stack = []
stack.append(x); x = stack.pop()   # LIFO

from collections import deque
q = deque()
q.append(x); x = q.popleft()       # FIFO / BFS
\end{lstlisting}

\begin{tbl}
\begin{tabularx}{\linewidth}{@{}R{.49\linewidth}L@{}}
\toprule
\textbf{Operation} & \textbf{Cost} \\
\midrule
\py{list.append(x)}, \py{list.pop()} & $O(1)$ amortized / $O(1)$ \\
\py{list.pop(0)} & $O(n)$ \\
\py{deque.append(x)}, \py{appendleft(x)} & $O(1)$ \\
\py{deque.pop()}, \py{popleft()} & $O(1)$ \\
\py{x in list} / \py{x in set/dict} & $O(n)$ / average $O(1)$ \\
list indexing / dict lookup & $O(1)$ / average $O(1)$ \\
heap push/pop / sorting & $O(\log n)$ / $O(n\log n)$ \\
\bottomrule
\end{tabularx}
\end{tbl}

Use \py{deque} for BFS queues, \py{heapq} for priority queues, \py{Counter} for frequencies, and \py{defaultdict(list)} for grouping or adjacency lists. Dicts and sets rely on hashing.

\subsection{EAFP: handle expected failures}

\begin{lstlisting}
try: value = data[key]
except KeyError: value = default
\end{lstlisting}

\section{20. Common Gotchas}

\begin{tbl}
\begin{tabularx}{\linewidth}{@{}R{.39\linewidth}L@{}}
\toprule
\textbf{Gotcha} & \textbf{Preferred practice} \\
\midrule
\py{xs.sort()} vs. \py{sorted(xs)} & mutates and returns \py{None} vs. returns new list \\
mutable function defaults & use \py{None}, then create inside \\
float equality & use \py{math.isclose(a,b)} \\
mutating during iteration & iterate over \py{xs.copy()} or build a new collection \\
\bottomrule
\end{tabularx}
\end{tbl}

\begin{lstlisting}
# BAD: rows refer to the same inner list
grid = [[0] * 3] * 3
# GOOD: a distinct list per row
grid = [[0] * 3 for _ in range(3)]
\end{lstlisting}

\begin{lstlisting}
# Filtering without modifying the traversed list
items = [x for x in items if keep(x)]
\end{lstlisting}

Avoid broad \py{except Exception} unless re-raising or at a deliberate application boundary.

\end{multicols}
\end{document}
